\section{Билет 10. Метод возможных направлений}

\ticketq{Метод возможных направлений.}

\paragraph{Постановка.}
\[
  \varphi_0(x) \to \min_{x \in S} ,
  \qquad
  S = \{\, x \in \R^n \mid \varphi_i(x) \le 0,\ i = \overline{1,m};\ Ax = b \,\} ,
\]
Метод строит допустимые точки $x_k \in S$ с убыванием цели по схеме
$x_{k+1} = x_k + \alpha_k s_k$, где $s_k$ — возможное направление спуска,
$\alpha_k > 0$ — шаг. Требуется допустимая начальная точка $x_0 \in S$.

\textbf{Предположения:} $\varphi_i$ выпуклы и дифференцируемы, градиенты
ограничены и липшицевы,
\[
  \norm{\nabla\varphi_i(x)} \le K , \qquad
  \norm{\nabla\varphi_i(x) - \nabla\varphi_i(y)} \le R\norm{x-y} ,
\]
$S$ компактно, выполнено условие невырожденности.

\paragraph{Определения.}
Ненулевой $s$ — \textbf{возможное направление} в $x \in S$, если
$\exists\,\delta > 0 : x + \alpha s \in S$ при всех $0 < \alpha < \delta$.
Если дополнительно $\varphi_0(x + \alpha s) < \varphi_0(x)$, то $s$ —
\textbf{возможное направление спуска}.

Ограничения в точке $x$: \textbf{активные}
$I_0(x) = \{i \mid \varphi_i(x) = 0\}$; \textbf{пассивные}
$\{i \mid \varphi_i(x) < 0\}$; \textbf{почти активные}
$I_\delta(x) = \{i \mid -\delta \le \varphi_i(x) \le 0\}$ — учитываются,
чтобы точка не подходила вплотную к границе.

\paragraph{Лемма (о возможном направлении).}
\textit{Пусть $x \in S$ и}
\[
  \nabla\T\varphi_i(x)\,s < 0 \ \ \forall i \in I_0(x) ,
  \qquad As = 0 .
\]
\textit{Тогда $s$ — возможное направление в $x$.}

\textsc{Док-во.} \emph{Равенства:} $Ax = b$, $As = 0$ $\Rightarrow$
$A(x + \alpha s) = b + \alpha\cdot 0 = b$ при любом $\alpha$.

\emph{Активные $i \in I_0(x)$:} $\varphi_i(x) = 0$, по формуле Тейлора
\[
  \varphi_i(x + \alpha s) = \alpha\,\nabla\T\varphi_i(x)s + o(\alpha) ,
\]
главный член отрицателен $\Rightarrow$ $\varphi_i(x+\alpha s) < 0$ при малом
$\alpha > 0$.

\emph{Пассивные:} $\varphi_i(x) < 0$ и $\varphi_i$ непрерывна $\Rightarrow$
$\varphi_i(x + \alpha s) < 0$ при малом $\alpha$.

Значит $x + \alpha s \in S$. $\blacksquare$

\paragraph{Лемма (о направлении спуска).}
\textit{$\nabla\T\varphi_0(x)s < 0$ $\Rightarrow$ $s$ — направление спуска.}

\textsc{Док-во.}
$\varphi_0(x+\alpha s) - \varphi_0(x) = \alpha\nabla\T\varphi_0(x)s + o(\alpha)$;
главный линейный член отрицателен $\Rightarrow$ при малом $\alpha > 0$ разность
отрицательна. $\blacksquare$

Итого: одновременное выполнение $\nabla\T\varphi_0(x)s < 0$,
$\nabla\T\varphi_i(x)s < 0$ ($i \in I_0$), $As = 0$ даёт возможное направление
спуска.

\paragraph{Вспомогательная задача Зойтендейка.}
Направление ищется как решение задачи ЛП: при $\xi_i > 0$, $i = \overline{0,m}$,
\[
  \eta \to \min_{\mathbb{H}} ,
  \quad
  \mathbb{H} = \Big\{ \eta \ \Big|\
    \nabla\T\varphi_i(x)s \le \xi_i\eta,\ i \in I_\delta(x)\cup\{0\} ;\
    As = 0 ;\ \norm{s}_\infty \le 1 \Big\} .
\]
Пара $s = 0$, $\eta = 0$ допустима $\Rightarrow$ $\eta_\delta \le 0$.

\textbf{Смысл.} При $\eta < 0$ и $\xi_i > 0$ имеем $\xi_i\eta < 0$, то есть
$\nabla\T\varphi_0 s < 0$ и $\nabla\T\varphi_i s < 0$ — по леммам $s$ есть
возможное направление спуска. Условие $As = 0$ сохраняет равенства, а
$\norm{s}_\infty \le 1$ нормирует задачу: без него $s$ можно было бы
растягивать и $\eta \to -\infty$. Случай $\eta_\delta = 0$ означает, что
возможного направления спуска нет, то есть точка приближённо стационарна.

\paragraph{Лемма (оценки вдоль направления).}
\textit{При $\alpha > 0$, $\norm{s_\delta} \le 1$:}
\[
  \varphi_0(x + \alpha s_\delta) \le \varphi_0(x) + \alpha\xi_0\eta_\delta
    + \alpha^2 R\norm{s_\delta}^2 ,
  \quad
  \varphi_i(x + \alpha s_\delta) \le \varphi_i(x)
    + \alpha\nabla\T\varphi_i(x)s_\delta + \alpha^2 R\norm{s_\delta}^2 ,
\]
\textit{и $A(x+\alpha s_\delta) = b$.} Первое — из липшицевости градиента и
ограничения $\nabla\T\varphi_0 s \le \xi_0\eta_\delta$.

\paragraph{Допустимость и убывание при $\eta_\delta < 0$.}
Разбор по трём группам даёт явные границы шага:
\begin{csteps}
  \item $i \in I_\delta(x)$ (почти активные): $\varphi_i(x) \le 0$ и
        $\nabla\T\varphi_i s_\delta \le \xi_i\eta_\delta < 0$, поэтому
        $\varphi_i(x+\alpha s_\delta) \le \alpha\xi_i\eta_\delta + \alpha^2 R
         \le 0$ при $\alpha \le \xi_i\abs{\eta_\delta}/R$;
  \item $i \notin I_\delta(x)$ (с запасом): $\varphi_i(x) < -\delta$ и
        $\norm{\nabla\varphi_i} \le K$, поэтому
        $\varphi_i(x+\alpha s_\delta) \le -\delta + \alpha K + \alpha^2 R \le 0$
        при малом $\alpha$;
  \item цель: $\varphi_0(x+\alpha s_\delta) - \varphi_0(x)
        \le \alpha\xi_0\eta_\delta + \alpha^2 R
        \le \tfrac12\alpha\xi_0\eta_\delta$ при
        $\alpha \le \xi_0\abs{\eta_\delta}/(2R)$.
\end{csteps}
Значит дробление $\alpha_k = \alpha_0\lambda^{\,j}$ конечно, шаг ограничен
снизу величиной, зависящей от $\abs{\eta_\delta}$, а $\varphi_0(x_k)$ строго
убывает. При $\eta_\delta = 0$ допустимого направления спуска нет —
выполняются условия оптимальности, и $\delta$ уменьшают. $\blacksquare$

\paragraph{Алгоритм.}
Задать $x_0 \in S$, $\eps > 0$, $\delta_0 > 0$, $\lambda \in (0,1)$.
\begin{csteps}
  \item Найти $I_0(x_k)$ и $I_{\delta_k}(x_k)$.
  \item Решить вспомогательную задачу; получить $s_k$, $\eta_k$.
  \item $\eta_k \ge 0$: направления нет — уменьшить $\delta_k := \lambda\delta_k$
        (при $\delta_k < \eps$ — конец, $x_k$ приближённо стационарна),
        к~шагу~1.
  \item $\eta_k < 0$: найти шаг дроблением — $\alpha_k = \alpha_0\lambda^{\,j}$
        с наименьшим $j$, при котором
        \[
          x_k + \alpha_k s_k \in S
          \quad\text{и}\quad
          \varphi_0(x_k + \alpha_k s_k) - \varphi_0(x_k)
            \le \tfrac12\,\alpha_k\,\xi_0\,\eta_k .
        \]
  \item $x_{k+1} = x_k + \alpha_k s_k$; останов при
        $\norm{x_{k+1}-x_k} < \eps$, иначе $k := k+1$, к~шагу~1.
\end{csteps}

\paragraph{Начальное приближение.}
Если $x_0$ неизвестна, решается вспомогательная задача
\[
  \eta \to \min , \qquad \varphi_i(x) - \eta \le 0,\ i = \overline{1,m},
  \qquad Ax = b .
\]
По условию Слейтера существует $\bar{x}$ с $\varphi_i(\bar{x}) < 0$,
$A\bar{x} = b$ $\Rightarrow$ $\min\eta < 0$, и найденная точка строго внутри
всех неравенств — ровно то, что нужно для старта. Это та же идея, что у
вспомогательной задачи симплекс-метода.

\paragraph{Ограничения.}
Нужна допустимая $x_0$; на каждой итерации решается задача ЛП; при сложных
нелинейных ограничениях трудно выбрать допустимый шаг; без учёта почти
активных ограничений метод неустойчив около границы; для невыпуклых задач —
только локальный минимум.

\seefull{гл.~3, \S1 и \S4 (численный пример в \S4)}
